\section{Introdução}

\vspace{-0.4cm}

\noindent\rule{\textwidth}{0.7pt}

Descreva aqui uma breve apresentação do capítulo 1.

\subsection{Motivação e Contexto}

<Análise e delimitação da situação inicial do problema de pesquisa. Aborde os assuntos relacionados ao tema que serão estudados. Detalhe as motivações por exemplo [científicas, técnicas, mercado, etc] relatando quais inovações e/ou pesquisas estão motivando o desenvolvimento do trabalho. Aqui você precisa despertar o interesse do leitor em sua pesquisa.>

\textcolor{red}{
<Os parágrafos produzidos precisam fazer uso de citações.
O recomendável é 1 citação (direta ou indireta)
para cada 2 parágrafos produzidos.>
}

\subsection{Problema de Pesquisa}

<Descreva com detalhes o seu problema de pesquisa.  <[Lembre-se do que foi conversado em sala de aula, sobre os elementos que compõem a problematização].>


\textcolor{red}{
<Os parágrafos produzidos precisam fazer uso de citações.
O recomendável é 1 citação (direta ou indireta)
para cada 2 parágrafos produzidos.>
}

<E por fim, resuma a problemática em uma pergunta
que irá nortear sua pesquisa.>

\subsection{Justificativa}

Descrever a justificativa para realização da pesquisa, explicando por que o problema é relevante. Aqui você precisa deixar clara a importância da sua pesquisa. Cite as vantagens que sua pesquisa irá trazer em relação ao problema de pesquisa.

\textcolor{red}{
<Os parágrafos produzidos precisam fazer uso de citações.
O recomendável é 1 citação (direta ou indireta)
para cada 2 parágrafos produzidos.> 
}

\subsection{Objetivos}
<Escrever, ao menos, um parágrafo introdutório para a seção de objetivos da pesquisa.>

\subsubsection{Objetivo Geral}

<Definição do objetivo geral.>

\subsubsection{Objetivos Específicos}

\vspace{0.7cm}

\begin{itemize}[label=$\bullet$]
    \item OE 1
    \item OE 2
    \item OE 3
    \item OE N
\end{itemize}